\PassOptionsToPackage{table,xcdraw}{xcolor}
\documentclass[10pt, a4paper]{article}

\usepackage[T1]{fontenc}
\usepackage[utf8]{inputenc}
\usepackage{tgpagella}
\usepackage{tgcursor}
\usepackage[spanish, es-noshorthands]{babel}
\usepackage{amsmath, amssymb}
\usepackage{booktabs}
\usepackage{tabularx}
\usepackage{xcolor}
\usepackage{tikz}
\usetikzlibrary{shapes.geometric, positioning}
\usepackage[american]{circuitikz}
\usepackage[most]{tcolorbox}
\usepackage[margin=1.6cm]{geometry}
\usepackage{fancyhdr}
\usepackage{lastpage}

\setlength{\headheight}{16pt}
\setlength{\emergencystretch}{2em}

\definecolor{ucbblue}{RGB}{0, 51, 102}
\definecolor{neutralbg}{RGB}{246, 248, 251}

\pagestyle{fancy}
\fancyhf{}
\lhead{\textbf{IMT-221 -- Circuitos Electrónicos II}}
\chead{Examen Escrito Individual EC2}
\rhead{23/11/2026}
\cfoot{\footnotesize Página \thepage\ de \pageref{LastPage} \quad|\quad Versión B -- Instrumento de 100 puntos}
\renewcommand{\headrulewidth}{0.5pt}
\renewcommand{\footrulewidth}{0.4pt}

\newcommand{\answerbox}[1]{%
    \vspace{0.15cm}
    \begin{tcolorbox}[colback=white, colframe=black!45, boxrule=0.5pt, arc=1mm, height=#1]
        \hfill {\tiny\color{black!45} Desarrollo, unidades e interpretación}
    \end{tcolorbox}
}

\begin{document}

% --- CABECERA PARA ESTUDIANTES ---
\begin{center}
    {\large\textbf{\textcolor{ucbblue}{Universidad Católica Boliviana ``San Pablo'' --- Sede Tarija}}}\\[0.08cm]
    {\Large\textbf{IMT-221 --- Circuitos Electrónicos II}}\\
    {\normalsize\textbf{Evaluación Escrita Individual EC2} \quad$\cdot$\quad 23 de noviembre de 2026 \quad$\cdot$\quad \textbf{Versión B}}
\end{center}

\begin{tcolorbox}[colback=neutralbg, colframe=ucbblue, boxrule=0.8pt, arc=1.5mm, left=6pt, right=6pt, top=5pt, bottom=5pt]
    \small
    \renewcommand{\arraystretch}{1.45}
    \begin{tabularx}{\textwidth}{@{}lX lX@{}}
        \textbf{Nombre:} & \hrulefill & \textbf{Versión:} & \textbf{B} \quad (Total: \textbf{100 puntos}) \\
        \textbf{Duración:} & \texttt{105 minutos} & \textbf{Calculadora:} & \texttt{Válida} \\
    \end{tabularx}
    \vspace{0.15cm}
    \hrule
    \vspace{0.15cm}
    \footnotesize
    \textbf{Instrucciones generales:} Responda de manera individual. Muestre todas las ecuaciones, sustituciones con unidades, supuestos e interpretación física. Los parámetros de los dispositivos son \textbf{modelos hipotéticos para este examen} y no corresponden a mediciones del curso ni a especificaciones de un fabricante particular. Durante la prueba no se realizan consultas conceptuales. No se solicitan ensayos físicos ni se acredita el PFI con esta evaluación.
\end{tcolorbox}

% =========================================================================
% PROBLEMA P1 (30 PTS VERSION B)
% =========================================================================
\section*{Problema P1 --- Amplificador operacional real (30 puntos)}

Un amplificador \textbf{no inversor} tiene $R_f = 100\,\mathrm{k\Omega}$ conectada entre la salida $v_o$ y la entrada inversora ($-$), y $R_g = 10\,\mathrm{k\Omega}$ conectada entre la entrada inversora ($-$) y la referencia de $0\,\mathrm{V}$. Considere un modelo de op-amp hipotético con $f_T = 1.1\,\mathrm{MHz}$, $\mathrm{SR} = 0.40\,\mathrm{V/\mu s}$, rieles de alimentación $\pm 9\,\mathrm{V}$ y excursión lineal de salida permitida hasta $\pm 7\,\mathrm{V}$. Se desea obtener una \textbf{salida sinusoidal} de $4.0\,\mathrm{V}$ de amplitud pico a una frecuencia de $20\,\mathrm{kHz}$. Para este problema suponga un modelo dominante de un polo y ganancia de ruido igual a la ganancia de señal ($\text{noise gain} = A_{CL}$).

\begin{center}
    \begin{circuitikz}[>=stealth, thick, scale=0.88, transform shape]
        \draw (0,0) node[op amp, yscale=-1] (op) {};
        \draw (op.+) to[short, -o] ++(-1.8,0) node[left]{$v_i$};
        \draw (op.out) to[short, *-o] ++(1.8,0) node[right]{$v_o$};
        \draw (op.-) -- ++(0,-1.1) coordinate(ninv);
        \draw (ninv) to[R, l=$R_g$] ++(0,-1.8) node[ground]{};
        \node[left, font=\footnotesize] at ([xshift=-0.35cm, yshift=-0.9cm]ninv) {$10\,\mathrm{k\Omega}$};
        \draw (ninv) to[short, *-] ++(0.6,0) to[R, l=$R_f$] ++(2.2,0) -| (op.out);
        \node[below, font=\footnotesize] at ([xshift=1.7cm, yshift=-0.35cm]ninv) {$100\,\mathrm{k\Omega}$};
        \draw (op.down) -- ++(0,0.6) node[above]{\footnotesize $+9\,\mathrm{V}$};
        \draw (op.up) -- ++(0,-0.6) node[below]{\footnotesize $-9\,\mathrm{V}$};
        \node[draw=ucbblue!60, fill=neutralbg, rounded corners, font=\scriptsize, align=left] at (5.6, -0.5) {%
            \textbf{Modelo hipotético:}\\
            $f_T = 1.1\,\mathrm{MHz}$\\
            $\mathrm{SR} = 0.40\,\mathrm{V/\mu s}$\\
            Rieles: $\pm 9\,\mathrm{V}$\\
            Excursión lineal máx.: $\pm 7\,\mathrm{V}$
        };
    \end{circuitikz}
\end{center}

\begin{enumerate}
    \renewcommand{\labelenumi}{(\alph{enumi})}
    \setlength{\itemsep}{2pt}
    \item Calcule la ganancia ideal en lazo cerrado $A_{CL}$ y el valor pico de la entrada $v_i$ necesario. \hfill \textbf{(6 pts)}
    \item Estime el ancho de banda de lazo cerrado $f_{BW}$ y determine si $20\,\mathrm{kHz}$ está por debajo de él. \hfill \textbf{(6 pts)}
    \item Calcule el \textit{slew rate} mínimo requerido en $\mathrm{V/\mu s}$ y compárelo con el disponible. \hfill \textbf{(8 pts)}
    \item Compruebe la excursión de salida y emita un dictamen conjunto de viabilidad. ¿Por qué superar el criterio de ancho de banda no basta? \hfill \textbf{(6 pts)}
    \item Calcule el pico máximo de salida permitido \textbf{solo por SR} para esa frecuencia e indique una limitación adicional que el modelo no descarta. \hfill \textbf{(4 pts)}
\end{enumerate}


% =========================================================================
% PROBLEMA P2 (30 PTS VERSION B)
% =========================================================================
\section*{Problema P2 --- Acondicionamiento diferencial / In-Amp y CMRR (30 puntos)}

Un acondicionador se modela en régimen lineal mediante:
\[
    v_o = A_d v_d + A_{cm} v_{cm} + V_{\mathrm{offset}}, \qquad v_d = v_+ - v_-, \qquad v_{cm} = \frac{v_+ + v_-}{2}.
\]
Los datos son \textbf{hipotéticos}: $v_+ = 1.012\,\mathrm{V}$, $v_- = 0.988\,\mathrm{V}$ (ambos respecto a la referencia de $0\,\mathrm{V}$), $A_d = 100\,\mathrm{V/V}$, $A_{cm} = 0.020\,\mathrm{V/V}$ y $V_{\mathrm{offset}} = +0.010\,\mathrm{V}$.

\begin{center}
    \begin{tikzpicture}[>=stealth, thick, scale=0.92, transform shape]
        \draw[fill=neutralbg, draw=ucbblue, rounded corners=2pt] (0,0) rectangle (4.2,2.4);
        \node[align=center, font=\small\bfseries, text=ucbblue] at (2.1, 1.55) {Bloque Acondicionador\\(In-Amp genérico)};
        \node[align=center, font=\scriptsize] at (2.1, 0.55) {$v_o = A_d v_d + A_{cm} v_{cm} + V_{\mathrm{offset}}$};
        \draw[->] (-1.8, 1.8) node[left]{$v_+ = 1.012\,\mathrm{V}$} -- (0, 1.8) node[right, font=\footnotesize]{$+$};
        \draw[->] (-1.8, 0.6) node[left]{$v_- = 0.988\,\mathrm{V}$} -- (0, 0.6) node[right, font=\footnotesize]{$-$};
        \draw[->] (4.2, 1.2) -- (6.0, 1.2) node[right]{$v_o$};
        \draw (2.1, 0) -- (2.1, -0.55) node[ground]{};
        \node[right, font=\scriptsize] at (2.25, -0.45) {Ref. $0\,\mathrm{V}$};
        \node[draw=black!40, fill=white, rounded corners, font=\scriptsize, align=left] at (8.5, 1.2) {%
            $A_d = 100\,\mathrm{V/V}$\\
            $A_{cm} = 0.020\,\mathrm{V/V}$\\
            $V_{\mathrm{offset}} = +0.010\,\mathrm{V}$
        };
    \end{tikzpicture}
\end{center}

\begin{enumerate}
    \renewcommand{\labelenumi}{(\alph{enumi})}
    \setlength{\itemsep}{2pt}
    \item Calcule $v_d$ en $\mathrm{mV}$ y $v_{cm}$ en $\mathrm{V}$. \hfill \textbf{(6 pts)}
    \item Determine por separado la contribución diferencial, la contribución de modo común y la de \textit{offset} en $v_o$; indique el valor total de $v_o$. \hfill \textbf{(8 pts)}
    \item Calcule el $\mathrm{CMRR}$ lineal y en decibelios ($\mathrm{dB}$). \hfill \textbf{(7 pts)}
    \item Si ambas entradas aumentan $+0.50\,\mathrm{V}$ manteniendo $v_d$ constante, estime el cambio en la salida $\Delta v_o$. \hfill \textbf{(5 pts)}
    \item Explique por qué el $\mathrm{CMRR}$ obtenido no puede presentarse como una medición física de los In-Amp construidos por el grupo. \hfill \textbf{(4 pts)}
\end{enumerate}

% =========================================================================
% PROBLEMA P3 (30 PTS VERSION B)
% =========================================================================
\section*{Problema P3 --- Diodo real y rectificación de precisión (30 puntos)}

Considere dos circuitos de \textbf{media onda positiva}, conectados separadamente a la misma fuente de señal senoidal y con cargas resistivas $R_L$ hacia $0\,\mathrm{V}$:

\begin{itemize}
    \setlength{\itemsep}{1pt}
    \item \textbf{Circuito D (convencional):} fuente $v_i \to$ ánodo (A) de diodo $D \to$ cátodo (K) conectado al nodo $v_{o,D} \to$ resistencia $R_L$ a tierra. Para este examen use un \textbf{umbral aproximado hipotético} $V_\gamma = 0.45\,\mathrm{V}$.
    \item \textbf{Circuito P (superdiodo no inversor idealizado):} $v_i$ en entrada ($+$) del op-amp; entrada ($-$) conectada al nodo rectificado $v_{o,P}$; salida interna $v_a$ conectada al ánodo (A) del diodo $D$, cuyo cátodo (K) es el nodo $v_{o,P}$; resistencia $R_L$ desde $v_{o,P}$ a tierra. La alimentación del op-amp se supone suficiente para el análisis idealizado de ambos semiciclos.
\end{itemize}

\begin{center}
    \begin{circuitikz}[>=stealth, thick, scale=0.85, transform shape]
        % Circuito D (Convencional)
        \node[font=\small\bfseries, text=ucbblue] at (1.8, 2.9) {Circuito D (Convencional)};
        \draw (0,0) node[ground]{} to[sV] (0,2.0) -- (0.8,2.0) to[D, l=$D$] (3.0,2.0) to[short, *-o] (3.8,2.0) node[right]{$v_{o,D}$};
        \node[left, font=\footnotesize] at (-0.5, 1.0) {$v_i(t)$};
        \node[below, font=\scriptsize] at (1.45, 1.85) {A};
        \node[below, font=\scriptsize] at (2.35, 1.85) {K};
        \draw (3.0,2.0) to[R, l=$R_L$] (3.0,0) node[ground]{};

        % Circuito P (Superdiodo)
        \node[font=\small\bfseries, text=ucbblue] at (9.6, 2.9) {Circuito P (Superdiodo no inversor)};
        \draw (7.8, 1.0) node[op amp, yscale=-1] (op) {};
        \draw (op.+) to[short, -o] ++(-1.2,0) node[left]{$v_i(t)$};
        \draw (op.out) to[short, *-] ++(0.6,0) coordinate(va) node[above]{\footnotesize $v_a$} to[D, l=$D$] ++(2.2,0) coordinate(vo) to[short, *-o] ++(0.8,0) node[right]{$v_{o,P}$};
        \node[below, font=\scriptsize] at ([xshift=0.65cm, yshift=-0.15cm]va) {A};
        \node[below, font=\scriptsize] at ([xshift=1.55cm, yshift=-0.15cm]va) {K};
        \draw (vo) to[R, l=$R_L$] ++(0,-1.8) node[ground]{};
        \draw (op.-) -- ++(0,-1.1) -| (vo);
    \end{circuitikz}
\end{center}

\begin{enumerate}
    \renewcommand{\labelenumi}{(\alph{enumi})}
    \setlength{\itemsep}{2pt}
    \item Para amplitudes \textbf{pico positivas} de entrada de $0.30\,\mathrm{V}$ y $1.00\,\mathrm{V}$, calcule los picos de salida del circuito D con el modelo dado. \hfill \textbf{(6 pts)}
    \item Calcule los picos de rectificación ideal $\max(v_i, 0)$ y los esperados del circuito P idealizado. \hfill \textbf{(6 pts)}
    \item Construya una tabla de estados para $v_i > 0$ y $v_i < 0$ del circuito P: estado del diodo, estado del lazo, $v_{o,P}$ y tendencia de $v_a$. \hfill \textbf{(10 pts)}
    \item Justifique por qué no puede usarse $v_+ \approx v_-$ durante todo el período. \hfill \textbf{(5 pts)}
    \item Indique \textbf{una limitación física} relevante del circuito P y explique por qué el valor $0.45\,\mathrm{V}$ no debe tratarse como una propiedad universal de todo diodo. \hfill \textbf{(3 pts)}
\end{enumerate}

% =========================================================================
% PROBLEMA P4-B (SOLO VERSION B)
% =========================================================================
\section*{Problema P4-B --- Interpretación básica de filtro pasa-bajos (10 puntos)}

Se proporciona la descripción teórica de \textbf{una sección} de filtro pasa-bajos de segundo orden, de ganancia en banda pasante $K = 1$, frecuencia natural $f_0 = 10\,\mathrm{Hz}$ y factor de calidad $Q = 0.5$. En una representación \textit{teórica ilustrativa}, la magnitud de la ganancia es aproximadamente $1\,\mathrm{V/V}$ a $1\,\mathrm{Hz}$ y \textbf{exactamente $0.50\,\mathrm{V/V}$ en $f_0 = 10\,\mathrm{Hz}$ según el modelo canónico}.

\begin{center}
    \begin{tikzpicture}[>=stealth, thick]
        \draw[->] (-2.0, 0) node[left]{$V_{\mathrm{in}}(j\omega)$} -- (0, 0);
        \draw[fill=neutralbg, draw=ucbblue, rounded corners=2pt] (0, -0.8) rectangle (6.2, 0.8);
        \node[align=center, font=\small] at (3.1, 0) {\textbf{Filtro pasa-bajos de 2.\textsuperscript{o} orden}\\$K = 1,\quad f_0 = 10\,\mathrm{Hz},\quad Q = 0.5$};
        \draw[->] (6.2, 0) -- (8.2, 0) node[right]{$V_{\mathrm{out}}(j\omega)$};
    \end{tikzpicture}
\end{center}

\begin{enumerate}
    \renewcommand{\labelenumi}{(\alph{enumi})}
    \setlength{\itemsep}{2pt}
    \item Identifique la banda pasante de baja frecuencia y la tendencia de la magnitud a frecuencias altas. \hfill \textbf{(3 pts)}
    \item Convierta la ganancia lineal de $0.50\,\mathrm{V/V}$ a decibelios ($\mathrm{dB}$). \hfill \textbf{(3 pts)}
    \item Explique por qué no se debe etiquetar automáticamente $f_0 = 10\,\mathrm{Hz}$ como la frecuencia de $-3\,\mathrm{dB}$, e indique qué información adicional requeriría para afirmar que se midió el desempeño físico del filtro. \hfill \textbf{(4 pts)}
\end{enumerate}

\end{document}
